\section{Farmer application and decision support}\label{sec:app}

\paragraph{Design constraints.}
The intended users have entry-level Android phones (about 4\,GB of RAM), limited data plans and varied literacy.
The application is a progressive web app written in plain JavaScript modules, with no framework, CDN or build step:
\begin{itemize}
  \item 69\,KB of JavaScript in total, of which about 29\,KB loads for the first screen; the chart and the other screens are loaded lazily;
  \item a service worker precaches the application shell and fetches API data network-first with a timeout;
  \item when offline, it falls back to cached responses and labels them with the time they were saved.
\end{itemize}
All 264 interface strings, including decision reasons, warnings, errors, accessibility labels and a ten-topic farmer guide, exist in English, Tamil, Telugu, Malayalam, Hindi and Kannada.
Automated tests enforce identical key sets, identical placeholders and native-script content in every language.
Commodity names are always shown in their original published form, with a localised label from a reviewed dictionary where one exists.
The translations are AI-drafted and \emph{provisional} until reviewed by native speakers with agricultural vocabulary.

\paragraph{Freshness and provenance.}
Every price shows its report date and age.
A banner appears whenever the dataset is more than three days old.
During development the deployment served an archive ending on 2026-04-21, and labelled it a historical snapshot. After the official back-fill it serves data to 2026-10-22, labelled as official daily reports. It never calls batch-refreshed data ``live''.
Forecast screens show the model version, the source and licence, the measured test error for that commodity and horizon, and the measured coverage of the interval.

\paragraph{Decision cards.}
The cards turn the 80\% interval $[\ell_h,u_h]$ (on the log-return scale) into a transparent rule.
The holding cost is $c_h = c_0 + w\cdot h/7$: a fixed 1\% plus 5\% per week for perishables, or 1\% per week for durables.
These costs are placeholders that the farmer can edit.
\begin{itemize}
  \item \textsc{hold} if $\ell_h\ge\ln(1+c_h)$;
  \item \textsc{sell} if $u_h\le 0$;
  \item \textsc{compare} if a market in the same district reported a price at least 5\% higher within three days;
  \item \textsc{watch} otherwise.
\end{itemize}
Every card is labelled a \emph{market outlook}, not a profit recommendation, and has a ``Why am I seeing this?'' panel listing the interval, costs, accuracy and freshness.
The client recomputes the card when the farmer enters transport or holding costs.
The rule's test-period hit rates are reported in Section~\ref{sec:decisions}.

\section{Threats to validity}\label{sec:threats}
\begin{itemize}
  \item \textbf{Internal.}
    Hyper-parameters, model family and conformal quantiles were all chosen on 2024 data.
    The search space was amended once, after a validation-only observation (Section~\ref{sec:protocol}).
    The panel was selected on 2021--2023 reporting density.
    For backtest folds inside that window, this selection uses within-fold information about which series report often, but no prices.
  \item \textbf{External.}
    Results describe well-reported mandis of 36 commodities.
    The test period includes a change in the publishing system, and the source is a third-party mirror of official data.
  \item \textbf{Construct.}
    Modal price is the published summary statistic, not the price an individual farmer received.
    Exactly-zero changes may partly reflect reporting practice rather than market stability.
  \item \textbf{Statistical.}
    Diebold--Mariano tests treat daily cross-sectional mean losses as one time series.
    We report one test per horizon and model against persistence, without a multiplicity correction; $p$-values are descriptive.
    The Newey--West lag equals the far edge of the target window ($h+3$; 7 for the next report). Doubling it to 20 at the 7-day horizon still gives $p\approx10^{-32}$.
\end{itemize}

\section{Ethical considerations}\label{sec:ethics}
\begin{itemize}
  \item \textbf{Harm from false certainty.} Forecast errors can hurt farmers who act on them. The interface therefore shows intervals, measured error and data age, never guarantees a price, and offers \textsc{watch} by default.
  \item \textbf{Privacy.} No personal data is collected, and there are no accounts or analytics.
  \item \textbf{Licensing.} The data is used under GODL-India with attribution.
  \item \textbf{Language.} Machine-drafted translations are marked provisional, because a mistranslated warning is itself a harm.
  \item \textbf{Evidence on information services.} Information services do not reliably raise farm-gate prices \cite{fafchamps2012sms,mitra2018asymmetric}. Welfare claims would require a field study, which we have not done.
\end{itemize}

\section{Reproducibility}\label{sec:repro}
All code, configuration and the experiment registry are in the repository.
The following are deterministic given the archive:
\begin{itemize}
  \item the silver, feature, model, prediction and gold layers;
  \item every table and figure in this paper.
\end{itemize}
These are produced by \texttt{scripts/run\_forecast\_pipeline.sh} and the analysis modules under \texttt{py/agri/analysis}.
Each experiment writes a manifest recording the git commit, the Spark and Java versions, the seed, the periods, row counts, the search space, the selected models and any warnings.
The test suites cover leakage, metric definitions, split purging, conformal calibration, the decision rule, batch/stream parity, de-duplication, late events, checkpoint recovery, the ingestion adapter, API contracts and internationalisation.
The raw archive is not redistributed and can be downloaded from its public source.
